% ============================================================================================
% Section 9 draft — CleanRL PPO + RND on Atari, train run 1.
% Spliced into RND_development_document.tex before \section{Distance to ground truth metrics}.
% Figure: code/2026-08-05_cleanrl-rnd-writeup/rnd_pipeline_comparison.pdf (standalone TikZ,
% compile with latexmk -pdf rnd_pipeline_comparison.tex in that folder).
% ============================================================================================
\clearpage
\section{CleanRL PPO \texorpdfstring{$+$}{+} RND on Atari}
\label{sec:cleanrl}

Every RND result so far in this document comes from one implementation --- the project's own, built
up over train runs 1 through 5 and reused unchanged by \Cref{sec:pm-am}. That implementation was
written against the original paper and the original repository, but it has never been checked against
a second, independently written RND. This section runs one: CleanRL's
\texttt{ppo\_\allowbreak rnd\_\allowbreak envpool.py}, a widely used single-file PPO $+$ RND on
Atari, at the upstream commit \texttt{fe8d8a0} (2026-04-20).

Two things make it worth the compute. First, it is a different RND on a different domain with a
different base learner --- pixels instead of coordinates, PPO instead of SAC, two value heads instead
of one reward stream --- so agreement or disagreement with our own RND is informative in a way that
another sweep of our own code is not. Second, its published result is a single seed. CleanRL reports
$7100$ on MontezumaRevenge against Burda et al.'s $8152$, and states plainly that one seed was all
the compute allowed. Thirty seeds turn that single number into a distribution.

\subsection{Train run 1}
\label{sec:cleanrl-trainrun1}

This run asks four questions.

\begin{enumerate}[label=\arabic*.,leftmargin=*]
\item \textbf{Does the published single-seed result hold up across seeds?} CleanRL's $7100$ is one
  sample from a distribution whose spread is unknown. Thirty seeds give a mean and a spread to
  compare against Burda et al.'s three-seed $8152$.
\item \textbf{What does the environment auto-reset defect cost?} CleanRL's own documentation records
  a bug in this file and says ``it does not seem to impact performance'', without evidence.
  \Cref{sec:cleanrl-bug} shows the defect is larger than the documentation implies --- it
  contaminates the intrinsic advantage of every preceding step in a rollout, not just the affected
  row --- and this run uses the corrected code.
\item \textbf{Where do the two RND implementations actually differ?}
  \Cref{tab:cleanrl-vs-project-rnd} and \Cref{fig:cleanrl-rnd-pipelines} answer this knob by knob.
  Several of the differences are ones this project deliberately chose against CleanRL, and having
  both stacks written down in one place is a prerequisite for attributing any result to any of them.
\item \textbf{What does a faithful reproduction of this workload cost on this cluster, and how fast
  can it be made to run?} The published run took about $250$ hours for one seed. The profiling in
  \Cref{sec:cleanrl-throughput} measures where that time goes and what is recoverable.
\end{enumerate}

% --------------------------------------------------------------------------------------------
% (1) The comparison table, then the comparison figure.
% --------------------------------------------------------------------------------------------
\paragraph{The two RND implementations, knob by knob.}
\Cref{tab:cleanrl-vs-project-rnd} lists every load-bearing design choice of the two stacks side by
side. The right-hand column is the \runref{run-5 original-small} RND, which is the RND of
\Cref{tab:pm-am-trainrun1-algos} and of every RND arm in \Cref{sec:pm-am}; its cell-by-cell
provenance is \Cref{tab:trainrun5-hparams}. Rows where the two agree are marked \emph{same}; the
remaining rows are the whole of the difference between the two implementations.

\begin{table}[H]
\centering
\footnotesize
\setlength{\tabcolsep}{3pt}
\renewcommand{\arraystretch}{1.2}
\resizebox{\textwidth}{!}{%
\begin{tabular}{@{}>{\raggedright\arraybackslash}p{3.4cm}
>{\raggedright\arraybackslash}p{5.6cm}
>{\raggedright\arraybackslash}p{5.6cm}@{}}
\toprule
\textbf{Knob} &
\textbf{CleanRL PPO $+$ RND} &
\textbf{This project's SAC $+$ RND} \\
& \small(\texttt{ppo\_\allowbreak rnd\_\allowbreak envpool.py} @ \texttt{fe8d8a0}) &
  \small(\runref{run-5 original-small}, \Cref{tab:pm-am-trainrun1-algos}) \\
\midrule
\multicolumn{3}{@{}l}{\textit{Base learner and environment}} \\
base RL algorithm & PPO, on-policy, $128$ parallel environment copies & SAC, off-policy, one
  environment copy, replay buffer $10^{6}$ \\
domain & Atari from pixels: \texttt{Montezuma\allowbreak Revenge-v5} via envpool & PointMaze and
  AntMaze from coordinates \\
observation the bonus sees & newest frame of a $4\times84\times84$ \texttt{uint8} stack, so
  $1\times84\times84$ & the full state vector: $4$-d PointMaze, $29$-d AntMaze \\
\midrule
\multicolumn{3}{@{}l}{\textit{Bonus input}} \\
what the bonus is computed on & next observation $s'$ & same \\
whitening & running mean/std per pixel, clip $[-5,5]$, denominator $\sqrt{\sigma^2}$ & running
  mean/std per dimension, clip $[-5,5]$, denominator $\sqrt{\sigma^2+10^{-8}}$ \\
warm-up before training & $50 \times 128 = 6400$ rollout steps of a random agent, on the
  \emph{training} environments & $6400$ environment steps of a random agent, on a \emph{fresh
  throwaway} environment so the training visit counts stay clean \\
\midrule
\multicolumn{3}{@{}l}{\textit{Networks}} \\
target network & $3$ convolutions ($32,64,64$) $\to$ flatten $\to$ linear $512$; frozen & linear
  $4\to256\to128$; frozen \\
predictor network & the same $3$ convolutions $\to$ linear $512 \to 512 \to 512$, i.e.\ \textbf{two}
  fully-connected blocks deeper than the target & the same trunk \textbf{one} $128$-wide block
  deeper: $4\to256\to128\to128$ \\
output width $m$ & $512$ & $128$ \\
activation & LeakyReLU on the convolutions, PyTorch default negative slope $0.01$; ReLU on the
  predictor's extra blocks & LeakyReLU throughout, negative slope $0.2$ (the TensorFlow default the
  original paper used) \\
weight initialization & orthogonal, gain $\sqrt2$ & same \\
bias initialization & zero & same \\
\midrule
\multicolumn{3}{@{}l}{\textit{Readout and predictor training}} \\
bonus readout & $b = \tfrac12\sum_{j=1}^{m} e_j^2$ (half sum over $512$ dimensions) & $b =
  \tfrac1m\sum_{j=1}^{m} e_j^2$ (mean over $128$ dimensions) \\
predictor loss & squared error, mean over the output width, then masked and averaged over the
  minibatch & squared error, mean over the output width, averaged over the replay batch \\
update proportion & $0.25$ --- a random quarter of each minibatch trains the predictor, following
  the paper's Table 5 & $1.0$ --- the whole batch trains the predictor, following the original
  repository code \\
predictor input at update time & the \emph{stored current} observation of the rollout, not the next
  observation the bonus was computed on & the next observation, the same input the bonus used \\
optimizer & Adam, learning rate $10^{-4}$, $\epsilon = 10^{-5}$, \textbf{shared with the policy};
  gradients clipped at global norm $0.5$ & Adam, learning rate $10^{-4}$, $\epsilon = 10^{-8}$,
  \textbf{its own optimizer}; no gradient clipping \\
\midrule
\multicolumn{3}{@{}l}{\textit{Intrinsic reward}} \\
reward normalization & divide by the running standard deviation $\sigma_F$ of a non-episodic forward
  filter $F_t = \gamma_{\mathrm{int}} F_{t-1} + b$, no mean subtracted & same mechanism, same
  $\gamma = 0.99$ \\
forward-filter discount & $\gamma_{\mathrm{int}} = 0.99$ & $0.99$ \\
how the bonus enters the objective & a \textbf{second value head} with its own non-episodic return
  at $\gamma = 0.99$; the two advantages are combined as $A = 1.0\,A^{\mathrm{int}} +
  2.0\,A^{\mathrm{ext}}$ & a \textbf{single reward stream}: the replay buffer stores $r^{\mathrm{ext}}
  + \beta\,r^{\mathrm{int}}$, so there is no intrinsic value head \\
episodic treatment of the bonus & the intrinsic return never resets at an episode boundary & the
  bonus is inside the one reward, so it inherits SAC's episodic treatment \\
bonus weight & fixed by $\mathtt{int\_coef}=1.0$ and $\mathtt{ext\_coef}=2.0$; not swept & $\beta$
  swept over $15$ values from $10^{-3}$ to $10^{4}$ \\
\midrule
\multicolumn{3}{@{}l}{\textit{Run scale}} \\
environment steps per run & $2\times10^{9}$ & $10^{6}$ \\
seeds & $1$ published; $30$ in this run & $100$ to $300$ per arm \\
discount on the extrinsic stream & $0.999$ & $0.99$ in \Cref{sec:pm-am} \\
\bottomrule
\end{tabular}%
}
\caption{The two RND implementations knob by knob. The left column is CleanRL's
\texttt{ppo\_\allowbreak rnd\_\allowbreak envpool.py} at commit \texttt{fe8d8a0}, read from the
source; the right column is the \runref{run-5 original-small} stack, whose per-cell provenance is
\Cref{tab:trainrun5-hparams}. ``same'' means the two implementations agree on that row.
\Cref{fig:cleanrl-rnd-pipelines} draws the same comparison as two pipelines. Four of these
differences are ones this project chose deliberately against CleanRL and recorded at the time: the
LeakyReLU slope of $0.2$ rather than $0.01$, the mean readout rather than the half sum, the update
proportion of $1.0$ rather than $0.25$, and feeding the predictor's update the next observation
rather than the stored current one.}
\label{tab:cleanrl-vs-project-rnd}
\end{table}

\Cref{fig:cleanrl-rnd-pipelines} draws the same two stacks as pipelines, read top to bottom, with the
rows aligned so that a horizontal line across the page compares the same stage in both. Stages the
two implementations share are gray; stages that differ are orange, with the difference written into
the box. The figure is drawn vertically and at the document's own text width, so its text prints at
full body size --- \Cref{fig:trainrun5-pipeline}, the horizontal pipeline of train run 5, is
$11.1$~inches wide and is therefore scaled to about half when included, which is why its labels are
hard to read.

\begin{figure}[H]
\centering
\includegraphics[width=\linewidth]{code/2026-08-05_cleanrl-rnd-writeup/rnd_pipeline_comparison.pdf}
\caption{The two RND pipelines side by side, each read top to bottom: CleanRL's PPO $+$ RND on the
left, the project's \runref{run-5 original-small} SAC $+$ RND on the right. The nine rows are
aligned, so each row compares the same pipeline stage in the two implementations. Gray marks a stage
that is the same in both; orange marks a stage that differs, with the difference stated inside the
box. The blue dashed arrow is the predictor gradient, which reaches the predictor only --- the target
is frozen in both. Generated by
\texttt{code/2026-08-05\_\allowbreak cleanrl-rnd-writeup/\allowbreak rnd\_\allowbreak
pipeline\_\allowbreak comparison.tex} (standalone TikZ; \texttt{latexmk -pdf} in that folder).}
\label{fig:cleanrl-rnd-pipelines}
\end{figure}

% --------------------------------------------------------------------------------------------
% (2) The environment table.
% --------------------------------------------------------------------------------------------
\paragraph{Environment specification.}
The environment is a single Atari game, \texttt{Montezuma\allowbreak Revenge-v5}, built by envpool
$0.6.6$ with four explicit keyword arguments and seventeen defaults. Because almost all of the
preprocessing lives inside envpool's C++ rather than in a stack of visible Python wrappers,
\Cref{tab:cleanrl-envspec} lists every configuration key that is in force and says where each value
comes from, so the environment can be rebuilt from the table alone.

\begin{table}[H]
\centering
\footnotesize
\setlength{\tabcolsep}{3pt}
\renewcommand{\arraystretch}{1.2}
\resizebox{\textwidth}{!}{%
\begin{tabular}{@{}>{\raggedright\arraybackslash}p{4.2cm}
>{\raggedright\arraybackslash}p{5.0cm}
>{\raggedright\arraybackslash}p{5.4cm}@{}}
\toprule
\textbf{Parameter} &
\textbf{Value} &
\textbf{Where the value comes from} \\
\midrule
\multicolumn{3}{@{}l}{\textit{Identity and spaces}} \\
environment id & \texttt{Montezuma\allowbreak Revenge-v5} & passed by CleanRL; the benchmark script's
  only environment \\
task and ROM & \texttt{montezuma\_\allowbreak revenge}, ROM shipped inside the envpool wheel &
  derived from the id; no download, so the environment builds on a node with no internet \\
observation space & \texttt{Box(0, 255, (4, 84, 84), uint8)} & envpool defaults
  \texttt{stack\_num=4}, \texttt{img\_height=\allowbreak img\_width=84},
  \texttt{gray\_scale=True} \\
channel order & oldest frame first, so index $3$ is the most recent frame & envpool layout; the RND
  bonus uses index $3$ only \\
action space & \texttt{Discrete(18)} & the ALE minimal action set, which for this game is the full
  set \\
parallel copies & $128$ & \texttt{num\_envs}, passed \\
\midrule
\multicolumn{3}{@{}l}{\textit{Preprocessing}} \\
frame skip & $4$ emulator frames per agent step, rewards summed & envpool default
  \texttt{frame\_skip=4} \\
max-pooling & element-wise maximum of the last two of the four skipped frames & envpool, fixed \\
grayscale and resize & ALE palette grayscale, then $210\times160 \to 84\times84$ with OpenCV
  \texttt{INTER\_AREA} & envpool defaults \texttt{gray\_scale=True},
  \texttt{use\_inter\_area\_resize=True} \\
no-ops at reset & uniform on $\{0,\dots,29\}$, then one FIRE frame & envpool default
  \texttt{noop\_max=30}; FIRE is unconditional in $0.6.6$ when FIRE is in the minimal action set \\
sticky actions & repeat-action probability $0.25$, applied by the ALE itself & passed by CleanRL,
  matching the original RND paper \\
\midrule
\multicolumn{3}{@{}l}{\textit{Episode}} \\
episodic life & \texttt{done} is raised on every life lost, not only on game over & passed
  (\texttt{episodic\_life=True}); the underlying game is not reset on a life loss \\
max episode steps & $108{,}000$ agent steps, so $432{,}000$ emulator frames & envpool's
  \texttt{registration.py} for every \texttt{-v5} id; in practice non-binding \\
termination rule & game over, or the step cap, or a life lost & envpool \texttt{atari\_env.h} \\
auto-reset & next-step: the terminating step returns the terminal observation, and the following
  step is consumed by the reset & envpool, fixed --- this is the defect
  \Cref{sec:cleanrl-bug} corrects \\
\midrule
\multicolumn{3}{@{}l}{\textit{Reward}} \\
reward seen by the agent & clipped to its sign, so $\{-1,0,+1\}$ & passed
  (\texttt{reward\_clip=True}) \\
reward recorded for scoring & the raw unclipped game score, kept in \texttt{info["reward"]} & envpool
  preserves it; this is what the episodic return in the run record accumulates \\
\bottomrule
\end{tabular}%
}
\caption{Environment specification for train run~1, in the style of \Cref{tab:pm-am-envspec}. Every
envpool configuration key in force is listed with its origin: ``passed'' means CleanRL supplies it
explicitly in \texttt{envpool.make}, everything else is an envpool default or registration value.
Two values are worth care and were verified against the installed $0.6.6$ rather than the web
documentation: \texttt{max\_\allowbreak episode\_\allowbreak steps} is $108{,}000$ \emph{agent}
steps (the web table's $27{,}000$ is a newer release's default), and the keys
\texttt{use\_\allowbreak fire\_\allowbreak reset} and \texttt{full\_\allowbreak action\_\allowbreak
space} do not exist in $0.6.6$ at all.}
\label{tab:cleanrl-envspec}
\end{table}

% --------------------------------------------------------------------------------------------
% (3) The run table.
% --------------------------------------------------------------------------------------------
\paragraph{Algorithms.}
\Cref{tab:cleanrl-trainrun1-algos} gives the run's hyperparameters in the style of
\Cref{tab:pm-am-trainrun1-algos}. Every algorithm value is CleanRL's default, unchanged; the run
adds no sweep. What the middle column records instead is the small number of things this run changes
about the \emph{code} rather than the algorithm, each of which is justified in the paragraphs that
follow the table.

\begin{table}[H]
\centering
\footnotesize
\setlength{\tabcolsep}{3pt}
\renewcommand{\arraystretch}{1.2}
\resizebox{\textwidth}{!}{%
\begin{tabular}{@{}>{\raggedright\arraybackslash}p{4.0cm}
>{\raggedright\arraybackslash}p{4.6cm}
>{\raggedright\arraybackslash}p{6.0cm}@{}}
\toprule
\textbf{Hyperparameter} &
\textbf{Value} &
\textbf{Source and note} \\
\midrule
\multicolumn{3}{@{}l}{\textit{RL algorithm}} \\
algorithm & PPO with two value heads & CleanRL default \\
policy and value network & $3$ convolutions ($32,64,64$) $\to$ $3136 \to 256 \to 448$, then a
  $448\to448$ extra block; actor $448\to448\to18$; two critics $448\to1$ & CleanRL default \\
parallel environments & $128$ & CleanRL default \\
rollout length & $128$ steps, so a batch of $16{,}384$ & CleanRL default \\
minibatches per epoch & $4$ & CleanRL default \\
update epochs & $4$ & CleanRL default \\
learning rate & $10^{-4}$, annealed linearly to $0$ & CleanRL default, shared by the policy and the
  RND predictor \\
optimizer & Adam, $\epsilon = 10^{-5}$ & CleanRL default \\
clip coefficient & $0.1$; value loss clipped & CleanRL default \\
entropy coefficient & $0.001$ & CleanRL default \\
value coefficient & $0.5$ & CleanRL default \\
gradient clipping & global norm $0.5$ & CleanRL default \\
advantage normalization & per minibatch & CleanRL default \\
extrinsic discount & $\gamma = 0.999$ & CleanRL default \\
GAE $\lambda$ & $0.95$ & CleanRL default \\
\midrule
\multicolumn{3}{@{}l}{\textit{Intrinsic bonus}} \\
bonus & Random Network Distillation & N/A --- the run has one arm \\
intrinsic discount & $\gamma_{\mathrm{int}} = 0.99$, non-episodic & CleanRL default \\
extrinsic and intrinsic weights & $\mathtt{ext\_coef} = 2.0$, $\mathtt{int\_coef} = 1.0$ & CleanRL
  default; no sweep \\
update proportion & $0.25$ & CleanRL default \\
observation-normalization warm-up & $50$ rollouts of a random agent, $6400$ steps & CleanRL default \\
RND stack & as the left column of \Cref{tab:cleanrl-vs-project-rnd} & read from the source \\
\midrule
\multicolumn{3}{@{}l}{\textit{Run bookkeeping}} \\
environment steps per run & $2\times10^{9}$ & CleanRL's benchmark value \\
seeds & $30$, \texttt{a\_seed} $= 1 \dots 30$ & this run; CleanRL published one seed \\
configs in this table & $1$(algorithm) $\times$ $1$(environment) $\times$ $30$(seed) $= 30$ &
  N/A --- no hyperparameter is swept \\
\midrule
\multicolumn{3}{@{}l}{\textit{Changes this run makes to the code, not to the algorithm}} \\
environment auto-reset defect & corrected & \Cref{sec:cleanrl-bug}; the only change that alters what
  the optimizer sees \\
metric logging & one JSON record per run, no Weights and Biases and no tensorboard & the record
  format of \Cref{sec:logging}, carrying the extrinsic and the intrinsic reward \\
checkpoint and resume & every $8$ hours, one checkpoint per run, previous one replaced &
  \Cref{sec:cleanrl-resume}; required because a run is far longer than the $4$-day partition limit \\
throughput options & the options that leave the computation unchanged are on; the options that
  change it are off & \Cref{sec:cleanrl-throughput} \\
\bottomrule
\end{tabular}%
}
\caption{Hyperparameters of train run~1, in the style of \Cref{tab:pm-am-trainrun1-algos}. ``N/A''
marks a knob that does not apply. Every algorithm value is CleanRL's own default and nothing is
swept, so unlike the sweep tables of \Cref{sec:pm-am} this table has a single configuration and its
config count is just the seed count. The last block lists the changes this run makes to the
\emph{code}; only the first of them changes what the optimizer sees.}
\label{tab:cleanrl-trainrun1-algos}
\end{table}

% --------------------------------------------------------------------------------------------
% (4) The bug.
% --------------------------------------------------------------------------------------------
\subsubsection{The environment auto-reset defect, and what it costs}
\label{sec:cleanrl-bug}

CleanRL's documentation page for this file carries a \texttt{bug} admonition. It states that envpool
and gym's vectorized environments disagree about which observation a terminating step returns, that
this makes the terminal state ``off by one'', and that ``it does not seem to impact performance, so
we take a note here and await for the upstream fix''. Two of those three claims need revising.

\paragraph{What envpool actually does.}
Define $s_{\mathrm{last}}$ as the final observation of an episode and $s_{\mathrm{new}}$ as the first
observation of the next one. envpool resets on the \emph{following} step: the step that ends an
episode returns $s_{\mathrm{last}}$ with \texttt{done=True}, and the next call to \texttt{step} is
consumed entirely by the reset, returning $s_{\mathrm{new}}$ with the submitted action discarded.
Measured directly on the installed envpool $0.6.6$ rather than taken from the documentation, on a
single CartPole copy where a reset state and a terminal state are told apart by their numbers alone:

\begin{quote}\ttfamily\footnotesize\raggedright
t=15\ \ done=True\ \ \ obs=TERMINAL\ \ \ \ elapsed\_step=16\\
t=16\ \ done=False\ \ obs=RESET-LIKE\ \ elapsed\_step=0
\end{quote}

and on \texttt{Montezuma\allowbreak Revenge-v5} with this run's exact options, where a lost life
raises \texttt{done} without resetting the game:

\begin{quote}\ttfamily\footnotesize\raggedright
t=457\ \ done=True\ \ \ lives=5\ \ elapsed\_step=458\\
t=458\ \ done=False\ \ lives=5\ \ elapsed\_step=458\ \ \ (the counter does not advance)
\end{quote}

A differential test settles that the discarded action really is discarded: two identically seeded
environment pools fed identical actions except on that one step stay byte-identical, while the same
injection one step earlier diverges immediately.

\paragraph{What that does to the rollout buffer.}
CleanRL stores \texttt{dones[step]} from the \emph{previous} iteration's \texttt{done}, so the
condition $\texttt{dones}[t]=1$ marks exactly the consumed steps. At such an index the buffer holds
an observation that is the old episode's terminal state, an action the environment threw away, a
reward that is always exactly zero, and an intrinsic bonus computed on the \emph{new} episode's first
frame. It is a transition that never happened, and it enters the policy gradient, both value losses,
the predictor's own loss, and the statistics that normalize the intrinsic reward. On
\texttt{Montezuma\allowbreak Revenge-v5} with \texttt{episodic\_life=True} it is about $1.2\%$ of
rows.

\paragraph{Why the damage is larger than that fraction suggests.}
The extrinsic stream repairs itself: at index $t-1$ the factor $1 - \texttt{dones}[t]$ is zero, so
the fabricated advantage cannot propagate backwards. The intrinsic stream has no such factor. The
intrinsic return is deliberately non-episodic --- Burda et al.\ make it so, and CleanRL follows them
by hardcoding the intrinsic terminal mask to $1.0$ --- and that is exactly what lets the fabricated
row leak. Perturbing only the fabricated row's intrinsic reward and re-running the file's own
generalized-advantage code moves the intrinsic advantage of \emph{every preceding row in the
rollout}, decaying as $(\gamma_{\mathrm{int}}\lambda)^{\Delta} \approx 0.94^{\Delta}$: full amplitude
at the row itself, about half ten rows earlier, still measurable at the head of a $128$-step rollout.
So a defect in $1.2\%$ of rows perturbs a large fraction of the intrinsic advantages, which is the
signal the whole method rests on. Whether that changes the final score is a separate question this
run does not answer; the claim here is only that ``it does not seem to impact performance'' was never
tested and is not obviously true.

\paragraph{The correction.}
Two edits, both keyed on $\texttt{dones}[t]$, which is the only reliable detector. The obvious
alternative --- reading envpool's \texttt{elapsed\_step} for a zero --- fails under
\texttt{episodic\_life=True}, because a lost life consumes a step without resetting the counter; on a
$600$-step probe it found $4$ of the $24$ consumed rows.

\begin{enumerate}[label=\arabic*.,leftmargin=*]
\item \textbf{Carry the advantage through the consumed row} so that row $t-1$ chains to row $t+1$ as
  though the row were absent. This stops the backward leak in the intrinsic stream.
\item \textbf{Drop the consumed rows from the update batch}, so no loss is ever evaluated on a
  transition that did not happen. The batch is then walked as exactly \texttt{num\_minibatches} equal
  slices with the remainder dropped; a ragged final slice of one row makes the advantage
  normalization divide by a zero standard deviation, and the resulting NaN reaches the weights within
  one optimizer step.
\end{enumerate}

Neither upgrading envpool nor switching to its Gymnasium-style interface avoids this. The behaviour
is unchanged in every release through $1.2.5$, and Gymnasium $1.0$ moved its own vector API
\emph{to} next-step auto-reset to match envpool --- so the ``upstream fix'' the documentation waits
for is not coming, because envpool's behaviour became the standard. The run also corrects a smaller
consequence of the same defect: the episode-length counter in CleanRL's statistics wrapper
incremented on consumed steps, overstating episode length by one per life lost.

% --------------------------------------------------------------------------------------------
% (5) Logging and resumability.
% --------------------------------------------------------------------------------------------
\subsubsection{Logging and resumability}
\label{sec:cleanrl-resume}

\paragraph{Logging.}
The upstream file writes to tensorboard and optionally to Weights and Biases. This run writes
neither. It writes one JSON record per run in the format of \Cref{sec:logging}: the configuration and
identity at the top level, then \texttt{train\_\allowbreak episode\_\allowbreak history} with one
entry per finished episode carrying the extrinsic return, the intrinsic bonus and the episode length;
\texttt{train\_\allowbreak history} with the per-update means of both rewards; and
\texttt{eval\_\allowbreak history} with throughput, the losses, and the count of rows dropped by the
auto-reset correction. The record is rewritten in full at every logging step through a temporary file
and an atomic rename, and carries \texttt{completed} so a partial record is never mistaken for a
finished run.

One deviation from \Cref{sec:logging} is forced by scale. A $2\times10^{9}$-step Atari run finishes
on the order of a million episodes, and an unbounded per-episode list would outgrow the filesystem.
The record therefore keeps every episode up to $50{,}000$ and every hundredth after that, and stores
both the cap and the stride alongside the true episode count, so an analysis can always tell a run
with few episodes from a run whose episodes were mostly not kept.

\paragraph{Why resumability is required, and what is saved.}
The \texttt{gpu} partition kills a job at $4$ days. At the throughput measured in
\Cref{sec:cleanrl-throughput} a single $2\times10^{9}$-step run needs far longer than that, so a run
that cannot resume cannot finish at all. PPO is on-policy, so there is no replay buffer to carry ---
the rollout buffer is rebuilt from scratch each iteration. What must be carried is the learned state,
the normalization statistics and the random streams:

\begin{itemize}[leftmargin=*]
\item the policy and both value heads, and the RND predictor;
\item \textbf{the RND target network}, which is frozen and never trained but is randomly initialized
  and therefore \emph{defines} the bonus --- a checkpoint that omitted it would look correct and
  silently restart exploration at every resume;
\item the Adam state over the policy and the predictor together;
\item the observation normalizer, the intrinsic-reward normalizer, and the forward filter behind it;
\item the step and update counters, which also carry the learning-rate schedule;
\item the Python, NumPy, and both Torch random-number states.
\end{itemize}

The checkpoint measures $50.2$~MB, so all $30$ seeds are checkpointed rather than a subset, and the
$30$ checkpoints together are $1.5$~GB. It is written every $8$ hours to a temporary file and renamed
over the previous one, so exactly one checkpoint per run exists at any moment and a kill during a
write cannot destroy the last good one.

\paragraph{What cannot be saved, stated plainly.}
envpool exposes no way to serialize its $128$ emulators, so a resumed run restarts its environments
from a fresh reset. The cost is bounded --- at most one partial episode per environment copy per
resume, against an $8$-hour checkpoint interval --- but it is a real gap and not a rounding error,
and it means a resumed run is not bit-identical to an uninterrupted one even though every learned
quantity is.

\paragraph{Verified, not assumed.}
The resume was demonstrated before any long run was submitted. A short run was stopped after $8$
updates at step $1024$ with \texttt{completed=false}; re-running the identical command logged
\texttt{[resume] continuing at update 9 of 20, global\_step 1024}, ran to update $20$, and finished
with \texttt{completed=true}, a history whose update column reads $5, 10, 15, 20$ with no duplicated
or missing rows, and an accumulated runtime spanning both segments.
